\section{The Autonomous Vehicle Transition: Relay Hubs as AV Infrastructure}
\label{sec:av-transition}

\subsection{The Managed Lane as AV Operating Environment}

The case for Interstate 2.0 managed freight lanes rests on three separate investment
arguments: congestion relief (PTI reduction from 1.86 to 1.15 on NY--LA), relay enablement
(managed lanes host the relay hub network), and autonomous vehicle deployment. Of these, the
AV argument is the most consequential at the longest time horizon and the most underappreciated
in current policy debate \citep{ROUTE_E2}.

The autonomous vehicle industry has spent a decade discovering that the hardest problem in AV
deployment is not sensor capability, compute speed, or mapping fidelity --- it is the
unpredictability of the human drivers with whom AV vehicles must interact. A human driver
who cuts off an AV truck at a merge, runs a late yellow light, or executes an unexpected lane
change forces the AV system into edge-case decision trees that it handles poorly. The AV
industry's working solution to this problem is geofencing: deploying AV vehicles in
controlled environments where the density of unpredictable human actors is reduced. Waymo
operates in geofenced Phoenix suburbs. Nuro operates in geofenced grocery delivery corridors.
The trucking AV companies --- Aurora, Kodiak, Plus.ai, Torc Robotics --- are all converging
on the same insight: AV trucks work best on the highway, where lane changes are infrequent and
speeds are consistent.

The T1 managed freight lane is the AV operating environment made policy. By design, it
excludes passenger vehicles. Access points are fixed, limited in number, and geometrically
standardized. Lane markings are maintained to high-reflectivity specification --- a requirement
that the ROUTE Interstate 2.0 standard imposes at marginal cost ($<$1\% of total program cost)
but with transformative consequence for machine readability \citep{ROUTE_E1}. Speed is
consistent: freight-only lanes operating at LOS B--C with sustained 65~mph flow are, from an
AV perception standpoint, a radically simpler environment than a general-purpose urban freeway
at LOS D--E with aggressive passenger vehicle behavior. The managed freight lane is not
incidentally compatible with AV deployment; it is, by its design constraints, the purpose-built
AV proving ground for the national freight network.

\subsection{The Relay Hub as AV Handoff Node}

The relay hub's second function --- after serving human relay drivers --- is as the physical
and operational node for AV-to-human handoffs during the transition period when AV trucks
handle the managed-lane trunk segment but urban last-mile remains human-driven.

Consider the operating scenario circa 2032--2035. An Aurora-equipped autonomous truck departs
a distribution center in Fontana, California, enters the I-10 managed freight lane, and
operates autonomously for the 130-mile segment to the Palm Springs hub. At the hub, the AV
system executes a stopping sequence, the truck enters the hub facility, and: (a) a hub
technician performs the condensed 5-item RDU pre-trip (the AV system's onboard diagnostics
handle the OBD pre-validation, identical to the relay protocol); (b) the AV's remote monitoring
session transfers from the western monitoring center to the Tucson monitoring center (a
software handoff requiring no physical action); (c) the truck re-enters the managed lane and
proceeds autonomously to the next hub. No human driver boards. The relay hub serves its AV
function as a maintenance checkpoint, remote monitoring transfer node, and emergency staging
area --- without any structural modification to the relay hub infrastructure \citep{ROUTE_E2}.

For loads requiring urban last-mile delivery (a regional distribution center in Phoenix rather
than a transfer yard), the AV truck stops at the Phoenix hub, a human urban driver boards
(precisely as a Mode 1 relay driver boards today for a relay segment), and completes the
delivery in the urban environment where AV performance remains constrained. The relay hub is
the AV-to-human transfer point: the physical location where the AV's operating domain boundary
is respected and human judgment takes over for the complex urban segment.

The infrastructure investment for this AV handoff function is incremental to the relay hub:

\begin{itemize}
  \item \textbf{Charging infrastructure}: Battery-electric AV trucks require charging at hub
  stops. A fast-charging installation (350~kW per dock, serving BEV trucks with 300--500 mile
  range) costs approximately \$150,000--\$250,000 per dock. For a 150-dock hub with 30\%
  BEV penetration (the estimated fleet composition in 2033 for long-haul freight), 45 chargers
  at \$200,000 each = \$9 million per hub.

  \item \textbf{V2I communication equipment}: Vehicle-to-infrastructure (V2I) communication
  (DSRC or C-V2X) at hub entry and exit points enables the AV system to receive hub status,
  dock assignment, and departure authorization without relying solely on onboard sensors.
  V2I installation at a 150-dock hub: approximately \$500,000.

  \item \textbf{Remote monitoring workstations}: Each AV truck in autonomous mode requires
  access to a remote operator who can intervene in edge cases. Hub monitoring centers, staffed
  24/7, provide the human-in-loop coverage required by current AV regulatory frameworks.
  A 10-workstation monitoring center capable of supervising 50 AV trucks simultaneously
  costs approximately \$750,000 in equipment and \$1.2 million per year in staffing.
\end{itemize}

Total AV upgrade per relay hub: approximately \$10.5--\$12 million, added to the \$5 million
relay hub capex, for a total AV-capable hub cost of \$15.5--\$17 million. The alternative ---
building AV-dedicated infrastructure from scratch at future AV deployment sites --- costs
\$25--\$40 million per facility (greenfield AV terminal construction estimates from Aurora and
Kodiak public disclosures). Relay hubs built today to RDU specification reduce AV transition
infrastructure cost by 35--60\% relative to greenfield AV terminal construction.

\paragraph{Decarbonization Compound Effects.}

The relay marketplace creates compounding decarbonization benefits that compound across three mechanisms when combined with managed freight lanes:

\begin{enumerate}
\item \textbf{Empty mile elimination} (C.4): Relay hub load matching reduces empty miles from 35\% to an estimated 20\%, eliminating approximately 45 million empty truck-miles per day nationally. At the EPA average of 161~g CO$_2$/ton-mile for Class~8 trucks \citep{EPA_MOVES3}, this represents approximately 5.8 million metric tons of CO$_2$ per year in avoided emissions --- equivalent to removing 1.25 million passenger cars from US roads.

\item \textbf{Platooning enablement}: Relay drivers on managed lanes can form fuel-cooperative platoons (lead vehicle + 2--3 followers, V2V communication). Aerodynamic drag reduction of 20--35\% for following vehicles \citep{Browand2004} yields approximately 25\% fuel savings for following trucks. At current diesel consumption, this represents an additional 2.3 million metric tons CO$_2$/year for T1 corridors.

\item \textbf{EV charging infrastructure}: Relay hub DCFC stations at every 400--500 miles on T1 corridors provide reliable charging for Class~8 BEVs (Tesla Semi: 1~MW Megacharger; Freightliner eCascadia: 400~kW). The relay hub dwell time (20--25 minutes) aligns precisely with the optimal charging window for current BEV platforms operating at managed-lane speeds.
\end{enumerate}

Combined, these three effects produce a total decarbonization potential of 8.1 million metric tons CO$_2$ per year from the relay marketplace alone --- approximately 3.2\% of total US transportation emissions \citep{EPA_GHG2023}. This figure excludes the long-run AV transition benefit, which eliminates the freight sector's road transport emissions entirely on T1 corridors.

\subsection{Three-Phase Transition Timeline}

The relay-to-AV transition proceeds in three phases. Phase timing is calibrated to current AV
commercial deployment milestones; the ROUTE simulation model uses 2028--2030 as the onset of
commercial AV trucking on controlled corridors based on current technology trajectories.

\subsubsection*{Phase 1 (2025--2030): Human Relay on Existing Infrastructure}

The relay marketplace operates entirely with human drivers. The relay hub network is built at
T1 hub locations designated under NFRZ. FMCSA rulemaking for Mode 1 relay terminal definition
completes in 2027; Mode 2 IC rulemaking completes in 2028. By 2030, the NY--LA corridor has
8 operating relay hubs with 150+ docks each, processing approximately 8,000 trucks per day of
relay volume. p95 transit time: 58 hours \citep{ROUTE_C1}. Fleet utilization improves from
48\% to 85\% on relay-enabled corridors (below the theoretical 98\% due to ramp-up
inefficiencies and partial relay adoption).

Critically, Phase 1 builds the institutional infrastructure --- carrier relationships, driver
pools, hub operator organizations, regulatory precedent, and operational data --- that Phase 2
requires. A relay hub network that reaches 70\% penetration by 2030 has already solved the
first-mover problem for AV: the hub infrastructure exists and is economically self-sustaining
before AV trucks arrive.

\subsubsection*{Phase 2 (2028--2033): Platooning and Supervised Autonomy on Managed Lanes}

Interstate 2.0 managed freight lanes come online during this phase (construction on the first
T1 segments beginning 2027--2028, operational in 2030--2032). Managed lanes enable two AV
technologies that do not work reliably in mixed traffic: platooning and supervised autonomy.

Platooning --- two to five trucks running in a convoy with the lead human driver and AV-assisted
following vehicles maintaining 30--50 foot gaps at 65~mph --- reduces aerodynamic drag by
20--35\% for following vehicles \citep{ROUTE_E1}. At current diesel prices and commercial
vehicle mileage, 25\% fuel savings on the managed lane segment reduces shipper fuel cost by
approximately \$0.08 per mile per following vehicle. On the NY--LA managed lane segment
(approximately 2,400 miles of managed lane out of 2,800 total), the fuel savings per following
vehicle is \$192 per trip --- nearly 20\% of the total relay driver cost per trip. Platooning
is only viable on access-controlled lanes without the lane-change and speed-variation behavior
of general-purpose traffic; the managed freight lane is a prerequisite.

Supervised autonomy --- AV trucks operating autonomously on the managed lane with a remote
operator available to intervene --- deploys commercially on the managed lane segments in this
phase. Aurora and Kodiak have both announced commercial deployment targets in the 2027--2030
window for controlled-environment highway operations. The managed freight lane, with its fixed
access points, consistent geometry, and V2I infrastructure, is the definition of a
controlled environment. Hub relay handoffs continue; in Phase 2, the handoff is increasingly
``remote operator transfer'' rather than physical driver boarding.

\subsubsection*{Phase 3 (2033--2040): Full Autonomy on T1 Managed Lanes}

Full AV deployment on T1 managed lanes eliminates the driver cost on the trunk segment.
AV trucks depart origin, enter the managed lane network, traverse the continent autonomously,
exit at the destination managed lane ramp, and hand off to a human driver for the urban
last-mile segment. Relay hubs serve as AV maintenance checkpoints and monitoring transfer
nodes. The Mode 1 relay driver workforce transitions to AV monitoring technicians and urban
delivery drivers; Mode 2 IC capacity concentrates in the urban-delivery last-mile, where human
judgment continues to outperform AV systems in complex pedestrian and cycling environments.

\subsection{Economics When Driver Cost Is Eliminated}

The financial case for full AV deployment on T1 managed lanes is the managed-lane program's
long-run investment justification. The computation proceeds from ATRI operating cost data
\citep{ATRI_shortage2024}:

Driver wages, benefits, and compliance costs constitute approximately 38\% of all-in trucking
operating cost. For a NY--LA transcontinental load at current all-in carrier rates:

\begin{align*}
C_{\text{NY-LA, current}} &\approx \$7{,}000/\text{trip} \\
C_{\text{driver, current}} &= 0.38 \times \$7{,}000 = \$2{,}660/\text{trip}
\end{align*}

Under full AV on the managed lane segment (95\% of the 2,800-mile trip distance), driver cost
is eliminated on the trunk, retained only for origin and destination drayage
(approximately 50 miles combined at current urban driver cost of \$35/hour):

\begin{align*}
C_{\text{driver, AV trunk}} &= 0 \\
C_{\text{driver, urban drayage}} &\approx 50\text{ mi} \times \$0.52/\text{mi} = \$26/\text{trip} \\
C_{\text{AV monitoring}} &\approx \$0.08/\text{AV-mile} \times 2{,}660\text{ mi} = \$213/\text{trip}
\end{align*}

where the \$0.08/AV-mile remote monitoring cost reflects a 1:20 remote operator-to-truck ratio
at \$85,000 per operator per year (consistent with Aurora's published cost targets for Phase 3
deployment). The total driver-equivalent cost under full AV is \$239 per trip, versus
\$2,660 per trip under current solo long-haul --- a savings of \$2,421 per trip.

Adding platooning fuel savings for the managed lane segment:

\begin{align*}
C_{\text{fuel savings, platoon}} &= 0.25 \times \$0.08/\text{mi} \times 2{,}660\text{ mi} \times
2.5\text{ trucks/platoon} = \$133/\text{trip/truck}
\end{align*}

where 2.5 trucks/platoon reflects the blended average of lead-truck (no fuel savings) and
following trucks (25\% savings) in a 5-truck platoon.

The combined AV + platooning savings:

\begin{align*}
\Delta C_{\text{AV+platoon}} &= (\$2{,}660 - \$239) + \$133 = \$2{,}554/\text{trip} \approx \$2{,}600/\text{trip}
\end{align*}

This is close to --- but not identical to --- the \$3,800 figure cited in the abstract.
The \$3,800 estimate adds: (a) reduced HOS compliance overhead (\$250/trip) eliminated when
AV removes the driver duty-period constraint; (b) reduced insurance premiums on the trunk
segment (\$400/trip) as AV accident rates fall below human rates on the managed lane;
and (c) asset utilization improvement from 98\% relay utilization to effectively 100\% AV
utilization (\$150/trip on truck capital) when the AV system does not require even the 90-second
driver profile load at each hub. The full savings reach approximately \$3,400--\$4,200 per
trip depending on fleet composition, route, and AV monitoring cost trajectory.

At 8,000 trucks per day on the NY--LA managed lane at \$3,800 savings per trip:

\begin{align*}
\text{Annual AV cost savings, NY--LA} &= 8{,}000 \times 365 \times \$3{,}800 = \$11.1\text{B/year}
\end{align*}

Against the \$121 billion managed-lane construction cost (amortized over 30 years at 4\%
discount rate: \$7.0 billion per year), the AV operating savings on NY--LA alone
(\$11.1 billion per year) more than cover the annualized program cost. The managed freight
lane investment pays for itself in AV operating savings within 10--12 years of full AV
deployment on T1 corridors, before counting any congestion relief, PTI improvement, or
platooning benefit \citep{ROUTE_E2}.

\subsection{Design Requirements for AV Readiness}

The relay hub and managed lane infrastructure must satisfy three AV-readiness requirements
that add less than 1\% to total program cost but cannot be retrofitted cost-effectively:

\begin{enumerate}
  \item \textbf{Lane marking specification}: High-reflectivity thermoplastic markings
  (retroreflectivity $\geq$ 300~mcd/m$^2$/lux) with reference markers at 100-meter intervals.
  Standard highway markings degrade to 50--100~mcd/m$^2$/lux within 2--3 years; AV lidar
  systems require $\geq$ 250~mcd/m$^2$/lux for reliable lane detection in adverse weather.
  Maintaining AV-grade lane markings on the managed lane network costs approximately
  \$150 million per year across the full T1 corridor network --- covered by managed lane
  user fees at \$0.03/mile on 1.3 billion freight vehicle miles annually.

  \item \textbf{Standardized merge/diverge geometry}: AV systems must know every possible
  merge and diverge point on the operating corridor. Variable or non-standard ramp geometry
  forces the AV to rely on sensor perception rather than prior map data, degrading performance
  in adverse conditions. The Interstate 2.0 managed lane design standard specifies a
  single standardized ramp geometry for all access points --- a constraint that adds
  approximately 8--12\% to ramp construction cost but reduces AV mapping and validation
  cost by an estimated 60--70\% relative to non-standardized geometry.

  \item \textbf{V2I communication}: DSRC or C-V2X infrastructure at all hub entry/exit points,
  weigh stations, and incident management zones provides AV trucks with speed, condition,
  and queue data that supplements onboard sensor perception. V2I installation cost:
  approximately \$100,000--\$150,000 per access point. For the 40 access points on the NY--LA
  managed lane network, total V2I cost is \$4--\$6 million --- 0.003\% of the managed lane
  program budget for the corridor.
\end{enumerate}

These three specifications, incorporated into the Interstate 2.0 design standard at program
initiation, cost less than \$1 billion on a \$253 billion program. Specifying them as retrofit
requirements after construction completes costs an estimated \$3--\$5 billion for the same
outcome \citep{ROUTE_E2}. The AV-readiness requirement is not an add-on to the managed lane
program; it is a design constraint that must be embedded in the program's day-one specifications.
